\chapter*{Prólogo}
\addcontentsline{toc}{chapter}{Prólogo}
\markboth{Prólogo}{Prólogo}

Este libro reúne, en formato de texto de estudio, el contenido del curso ICM3323 --- \emph{Elementos Finitos en Diseño Mecánico}. El objetivo no es reemplazar las clases ni las diapositivas de la cátedra, sino ofrecer una versión desarrollada en prosa, con las derivaciones completas y comentadas, de modo que sirva como material de repaso y consulta a lo largo del semestre.

El libro está organizado en partes que agrupan capítulos temáticamente afines. La \textbf{Parte I} cubre los fundamentos de mecánica del medio continuo necesarios antes de introducir cualquier discretización: el problema de valores de frontera de un sólido elástico y la teoría constitutiva que relaciona tensiones y deformaciones. La \textbf{Parte II} desarrolla, a partir de esos fundamentos, la formulación completa del método de los elementos finitos para el caso más simple e ilustrativo posible --- el elemento de barra --- cubriendo desde el análisis matricial clásico hasta la formulación isoparamétrica y la integración numérica de Gauss. Está pensado como un documento vivo: a medida que el curso avance, se irán agregando nuevas partes (elementos de viga, elasticidad 2D y 3D, dinámica) siguiendo la misma estructura y convenciones visuales.

\section*{Cómo usar este libro}

A lo largo del texto se utilizan cinco tipos de recuadros de color para distinguir el tipo de contenido:

\begin{itemize}[leftmargin=1.4em]
  \item[] \textcolor{ucred}{\textbf{Ecuación importante}} (marco granate): resultados centrales que conviene memorizar o tener siempre a mano --- típicamente la fórmula final de una derivación. Todas quedan además recopiladas en el \hyperref[anexo:formulario]{formulario del Anexo A}.
  \item[] \textcolor{ucblue}{\textbf{Definición}} (marco azul pizarra): introducción formal de un concepto o cantidad nueva.
  \item[] \textcolor{ucteal}{\textbf{Ejemplo resuelto}} (marco verde azulado): aplicación completa de la teoría a un caso concreto, tal como se resolvió en clases.
  \item[] \textcolor{ucgray}{\textbf{Teorema}} (marco gris, título en versalitas): derivaciones formales completas, con su demostración.
  \item[] \textcolor{ucamber}{\textbf{Observación}} (marco ámbar, sin numerar): comentarios, aclaraciones o advertencias puntuales que no forman parte del hilo principal de la derivación.
\end{itemize}

Toda ecuación importante y todo ejemplo, definición o teorema está numerado por capítulo y referenciado mediante hipervínculos internos del PDF, de modo que se puede navegar el documento haciendo clic en cualquier referencia cruzada (por ejemplo, ``ver \textcolor{ucblue}{Ecuación importante 4.2}'').
